
\textbf{Further Remarks on Algorithm~\ref{alg:algname}.\ } We use learning rate scaling to match the Adam update RMS norm \citep{liu_muon_2025} for Muon as well as \algname{}. This corresponds to a learning rate adjustment by a factor of $0.2 \sqrt{\max(m, n)}$ for a layer of shape $m \times n$. Moreover, \algname{} as outlined in Algorithm~\ref{alg:algname} uses simplified Nesterov momentum \citep{MethodSolvingConvex1983nesterov, bengio_2013_advances}, matching the PyTorch implementation of Muon. 

To keep Algorithm~\ref{alg:algname} concise, we describe the use of decoupled weight decay with \algname{} separately here. When decoupled weight decay with coefficient $\lambda$ is enabled, we replace the last line of Algorithm~\ref{alg:algname} with $\W \gets \mathrm{Diag}(\frac{\g}{\rbf}) \Rbf- \eta_t\lambda \W$ and refresh $\g \gets \RowNorm{\W}$ to preserve the invariant $\g = \RowNorm{\W}$. This is only one of several natural placements of weight decay under the $(\g, \Rbf)$-parameterization. For instance, one could equally well decay only the direction component $\Rbf$. The variant adopted above stays closest to the canonical decoupled formulation \citep{loshchilov_decoupled_2019} on the effective weight $\W$ and is the one used throughout our experiments. A systematic study of other alternatives is left to future work.
 
\textbf{General Setup.\ } If not noted otherwise, we run all experiments on top of the implementation of \cite{ajroldi2024plainlm}. We make our code available on github.\footnote{\url{https://github.com/kcc-lion/muown}} We train a transformer with Rotational Positional Embeddings \citep{su_rope}, RMSNorm \citep{zhang_rmsnorm}, SwiGLU MLPs \citep{shazeer_glu_2020}, expansion ratio of 8/3, and weight tying. We employ a warmup-stable-decay schedule for 2\% of training as warmup and 20\% as cooldown. In all but the distributed benchmarking experiment in Table~\ref{tab:timing-overhead}, we make use of the PyTorch implementation of Muon. We justify this deviation below. We employ batch size of 524k tokens per step training, training for a total token budget of 20$\times$ tokens/parameter \citep{hoffmann_training_2022}, if not otherwise noted. Regarding sequence length and training budget we make the following size-dependent choices:
\begin{itemize}
    % 20270207_wd_ablation
    \item For models of the 124M class, we choose 12 layers with hidden dimension of 768 and 12 attention heads. We train with sequence length 1024 for 5B tokens and validate on 4M tokens. We remark that due to weight tying the actual number of parameters is 124M.
    % 20260207_500M_optimizer_sweep for tuning, for spectral 20260221_500M_max_logit
    \item For models with 500M parameters, we choose 22 layers with hidden dimension of 1280 and 20 attention heads. We train with sequence length 2048 for 10B tokens and validate on 10M tokens. 
    % 20260221_1B_multinode
    \item For models with 1B parameters, we choose 24 layers with hidden dimension of 1792 and 28 heads. We train for 20B tokens with sequence length 2048, validating on 10M tokens. 
    \item For models with 2.7B parameters, we choose 32 layers with dimension 2560 and 40 heads and train for 53B tokens, keeping everything else as in the 1B setup. 
\end{itemize}

All models are trained on the 100B token sample of FineWeb-Edu \citep{lozhkov2024fineweb-edu} and tokenized with the GPT-NeoX-20B tokenizer \citep{black_gpt-neox-20b_2022}, if not noted otherwise. Experiments are performed on either of NVIDIA GH200 and NVIDIA H100 GPUs and we estimate the total GPU hours spent as roughly $\approx$18k hours, including experimentation, development, and final runs.


\textbf{Spectral Analysis (Fig.~\ref{fig:sv_histograms}, \ref{fig:spectral-norm-detailed}, \ref{fig:row-norm-detailed}, \ref{fig:lambda-max-detailed}, \ref{fig:spectral-histogram-detailed}, \ref{fig:spectral-distribution-lineplot}). } For the spectral analysis plots, we adopt the 500M setup outlined above and tune Muon and \algname{} with learning rate $\eta \in \{1 \times 10^{-3}, 2 \times 10^{-3}, 4 \times 10^{-3}\}$ and weight decay $\lambda \in \{ 0, 0.1 \}$. We report results for the best runs in terms of validation loss, i.e. $\eta = 2 \times 10^{-3}$ with and without weight decay for Muon and $\eta = 4 \times 10^{-3}$ without weight decay for \algname{}. For the setup with fixed row magnitude, we re-use the best config from \algname{}.

\textbf{Learning Rate Ablation (Fig.~\ref{fig:lr-ablation}).} For the ablation study in Fig.~\ref{fig:lr-ablation}, we use a grid of $\eta \in \{2.5 \times 10^{-4}, 5 \times 10^{-4}, 7.5 \times 10^{-4}, 10^{-3}, 2 \times 10^{-3}, 4 \times 10^{-3}, 6 \times 10^{-3}, 8 \times 10^{-3}\}$ for all optimizers apart from Lion. For Muon's weight decay parameter, we use $\lambda=0.1$, which we found to be optimal across a grid of $\lambda \in \{0.01, 0.03, 0.1, 0.3\}$ (see Fig.~\ref{fig:weight-decay-ablation}). Similarly, based on the results of Fig.~\ref{fig:weight-decay-ablation}, we choose $\lambda=0.01$ for \algname{} when using weight decay, despite having only a marginal effect. For Lion, we use $\eta \in \{ 8 \times 10^{-5}, 10^{-4}, 2.5 \times 10^{-4}, 5 \times 10^{-4}, 7.5 \times 10^{-4}\}$, as Lion typically has a larger update norm than other optimizers, requiring a lower learning rate \citep{chen_symbolic_2023}. This matches our observation of divergence for learning rates larger than these. Moreover, \cite{chen_symbolic_2023} prescribe a larger value for weight decay in the range of 3-10$\times$ of the one used for AdamW. We opt for $5\times$ the optimal value of Muon, yielding $\lambda=0.5$ for Lion. We retain the default values of $(\beta_1, \beta_2) = (0.9, 0.99)$ for Lion. For SOAP \citep{vyas_soap_2024}, we use the default hyperparameters from the original implementation: $(\beta_1, \beta_2) = (0.9, 0.95)$, $\epsilon = 10^{-8}$, a preconditioner update frequency of 10 steps, and bias correction enabled. We set the maximum precondition dimension to 10{,}000, which excludes the embedding and head layers (vocabulary size 50{,}280) from preconditioning while covering all other layers. Weight decay for SOAP is set at $\lambda=10^{-4}$, following \cite{vyas_soap_2024}.

\textbf{2.7B Runs (Fig.~\ref{fig:3B-plot}, Table~\ref{tab:1B-results}).\ } We adopt the 2.7B setup above and sweep $\eta \in \{7.5\times 10^{-4}, 10^{-3}, 2\times 10^{-3}\}$, with $\lambda=0$ for \algname{} and $\lambda \in \{0, 0.1\}$ for Muon. The $\eta=2\times 10^{-3}$ Muon runs were aborted after 12h on 16 GPUs as the validation loss was diverging.

\textbf{Weight Decay Ablation (Fig.~\ref{fig:weight-decay-ablation}).} We adopt the 124M setup described above and train for 5B tokens ($2\times$ Chinchilla-optimal), sweeping a two-dimensional grid over the learning rate $\eta \in \{7.5 \times 10^{-4}, 10^{-3}, 2 \times 10^{-3}, 4 \times 10^{-3}, 6 \times 10^{-3}\}$ and weight decay $\lambda \in \{0, 0.01, 0.03, 0.1, 0.3\}$ for both Muon and \algname{}. We report perplexity at $2\times$ the Chinchilla-optimal token count (i.e. $5$B tokens).

\textbf{Learning Rate and Width Ablation (Fig.~\ref{fig:width-ablation}).} We use the 124M architecture outlined above and vary the width as 512, 768, 1280, resulting in models of size 67M, 124M, and 308M. We train each size for a Chinchilla-optimal amount of tokens and use a logarithmically-spaced learning rate grid from $2^{-16}$ to $2^{-5}$. 

\textbf{Gradient Conditioning (Fig.\ \ref{fig:gradient-conditioning}).} The effective rank of a matrix is defined as $\mathrm{erank}(\boldsymbol{\sigma}) = \exp (-\sum_{i=1}^{\min(m, n)} p_i \log p_i)$ where $p_i = \frac{\boldsymbol{\sigma}_i}{\| \boldsymbol{\sigma}\|_1}$ and $\boldsymbol{\sigma} \in \mathbb{R}^{\min(m, n)}$ contains the singular values of the gradient matrix. We normalize the effective rank by dividing it with $\min(m, n)$. For Muon, we show the gradient with respect to $\W$, and for \algname{} the gradient with respect to $\Rbf$. Each line shows one of the 12 layers and the bold line the average. 

\textbf{Gradient Noise (Fig.\ \ref{fig:scaled_cm_comparison}).} We base the gradient noise analysis on our 124M weight decay ablation setup and consider the configuration with the best final validation loss. This results in Muon ($\eta = 4 \times 10^{-3}, \lambda = 0.1$) and \algname{} without weight decay ($\eta = 6 \times 10^{-3}, \lambda = 0$) run. For each of these runs we save the model and optimizer states after $t = 1900, 3800, 5700, 7600$ and the final $T = 9538$ iterations. For every checkpoint, we first compute the true gradient $\nabla_\W \mathcal L \pare{\Wt}$ for Muon, and $\nabla_\g \Ft\pare{\gt, \Rt}, \nabla_\Rbf \Ft \pare{\gt, \Rt}$ for \algname{} checkpoints over the whole $5$B token training dataset $\{\xi_1, \dots, \xi_T\}$. In a second pass we then calculate, for each weight matrix $\W_t^{(i)}$, resp. $\g_t^{(i)}, \Rt^{(i)}$,
\begin{equation*}
    \sigma_{\W_t^{(i)}}^2 = \frac 1 T \sum_{\tau = 1}^{T}
    \trnorm{\nabla_\W \mathcal L \pare{\Wt} - \nabla_\W \mathcal L \pare{\Wt, \xi_\tau}}^2,
\end{equation*}
and similarly $\sigma_{\Rt^{(i)}}$ and $\sigma_{\gt^{(i)}}$. The values of $\zeta_\W, \zeta_\g$ and $\zeta_\Rbf$ are known analytically to be $\zeta_\W = \zeta_\Rbf = \sqrt{\min\{m, n\}}$ for $\W, \Rbf \in \mathbb R^{m \times n}$ and $\zeta_\g = \sqrt{m}$ for $\g \in \mathbb R^m$. We compute the noise term separately for each weight matrix, plot the median across weights as a bold line, and show the interquartile range as a shaded region. We find that the noise term of \algname{} remains consistently below that of Muon, with the gap widening over the course of training.

\textbf{nanoGPT Speedrunning (Fig.~\ref{fig:modded-nanogpt-learning-rate}). } For the nanoGPT speedrunning \citep{modded_nanogpt_2024} experiments, we leverage the Muon record available on github.\footnote{\url{https://github.com/KellerJordan/modded-nanogpt/blob/2c7027462732f2f7a76fb4b021066442c7b29298/records/track_1_short/2024-10-10_Muon/train_gpt2.py
}} 

\textbf{Multi-Seed Run for 500M (Table~\ref{tab:500M-multiseed}).\ } We adopt the 500M setup above and report mean and standard deviation across three seeds for each configuration. The grid is $\eta \in \{1, 2, 4\}\times 10^{-3}$ for all optimizers. \algname{} is run with $\lambda=0$; for Muon we report both $\lambda=0$ and $\lambda=0.1$, the latter identified as optimal in the 124M weight-decay ablation (Fig.~\ref{fig:weight-decay-ablation}). SOAP follows its hyperparameters from the 124M learning-rate ablation above with $\lambda=10^{-4}$.


\textbf{Resource Overhead (Table~\ref{tab:timing-overhead}).} We compare the overall step times for Muon and \algname{} by running 1K steps for the 500M model outlined above on 4 GH200 devices. Here, we report the median runtime for Muon not based on the PyTorch implementation, but based on a custom version that mirrors the implementation details of \algname{} as closely as possible. The reason for this deviation is that the PyTorch implementation does not provide a version which distributes the optimizer step across ranks, making a fair comparison in this setup infeasible.

\textbf{Qwen2-0.5B (Table~\ref{tab:qwen2-0.5B}).\ } We use the model configuration available on HuggingFace\footnote{\url{https://huggingface.co/Qwen/Qwen2-0.5B/blob/main/config.json}}, but lower the base frequency of RoPE to 10,000 to align the configuration with the initial phase of pre-training as outlined in \cite[Section~3.2]{yang_qwen2_2024}. We use the 10B split of FineWeb-Edu tokenized with the Qwen2-Tokenizer. We use a grid of $\eta \in \{10^{-3}, 2 \times 10^{-3}, 4 \times 10^{-3}, 8 \times 10^{-3}, 10^{-2}\}$ and tune $\lambda \in \{0, 0.1\}$ for Muon. 

\textbf{Magnitude Optimizer Ablation (Table~\ref{tab:magnitude-optimizer}).\ } We adopt the 500M setup and compare three choices for the magnitude optimizer in Algorithm~\ref{alg:algname} across $\eta \in \{1, 2, 4\}\times 10^{-3}$, all run with $\lambda=0$: \emph{Fixed} freezes $\g_t = \g_1$ at initialization, exposing the reparameterization without optimizing it; \emph{Signum} replaces Adam with signSGD with momentum, and \emph{Adam} is the default of Algorithm~\ref{alg:algname}.

\textbf{Licenses.} The FineWeb-Edu dataset \citep{lozhkov2024fineweb-edu} is released under the Open Data Commons Attribution License (ODC-By). \texttt{plainLM} \citep{ajroldi2024plainlm} and \texttt{modded-nanoGPT} \citep{modded_nanogpt_2024} are released under the MIT License. Qwen2 \citep{yang_qwen2_2024} is released under the Apache 2.0 License.